{\selectlanguage{english}
\begin{abstract}
This dissertation concerns the theoretical studies of the influence of different types of band structure engineering on the electronic properties of group IV semiconducting materials. The aim is to find such systems that allow for optoelectronics applications.

Four methods were presented to manipulate the parameters of the material, along with a step-by-step procedure to achieve and process the results of calculations to acquire the necessary properties. 
For some specifications, a more detailed explanation was needed, due to difficulties in interpretation or the way of calculating it.
In case of each method: strains, alloying, structural changes, and quantum confinement, the most important physical effects are shown along with the explanation of the mechanism.
In materials where the electronic properties allow us to use them in optoelectronic systems, the most important parameters for applications, e.g. electrons and holes mobility, are estimated.

In addition to presenting results and their analysis, the dissertation contains an expanded theoretical introduction related to all the problems discussed.
All useful derivations and proofs can be found in \hyperref[part:IV]{Appendix} attached to the thesis.
\end{abstract}

\ornament


\noindent\small
\textbf{Key words} \quad DFT \quad electronic structure \quad group IV semiconductors \quad strain \quad\\ alloy \quad quantum confinement
}

\vfil\vfil






\cleardoublepage